\section{RESUMEN}
El phishing sigue creciendo porque se centra en atacar a la persona, no al sistema, y las listas negras siempre llegan tarde, ya que cuando un sitio nuevo aparece reportado, ya cumplió su objetivo. Frente a esto, el proyecto propone un agente de IA que reciba una URL y estime su riesgo de phishing sin depender de esas listas, analizando en su lugar rasgos de la URL y el dominio (longitud, caracteres raros, IP en vez de dominio, edad de registro) y, cuando el sitio sea accesible, elementos del contenido como formularios de login o similitud con marcas conocidas. Para el modelo se utilizará XGBoost, Random Forest y LightGBM con el fin de promediar la importancia que cada uno asigna a las variables y quedarse con las más relevantes, pues sobre ese conjunto reducido se construirá un stacking de dos niveles, donde XGBoost y Random Forest generan cada uno una probabilidad de phishing que alimenta a un segundo XGBoost más simple encargado de la predicción final. Todo se entrenará con URLs reales de sitios de phishing y de sitios legítimos recolectadas de fuentes públicas, de modo que el agente reciba una URL cruda y no características ya calculadas. La meta es entregar un agente que funcione y claridad sobre los límites referente a qué tanto generaliza a sitios nuevos y qué tan fácil es engañarlo con URLs hechas para evadirlo.
\newline \par
\textbf{Palabras Clave:} Phishing, Ciberseguridad, Inteligencia Artificial, Aprendizaje Automático, Detección de URL, Análisis de Riesgo, Ensambles, Explicabilidad.
    